\documentclass[10pt]{article}

\usepackage[utf8]{inputenc}

\usepackage[T1]{fontenc}

\usepackage{amsmath}

\usepackage{amsfonts}

\usepackage{amssymb}

\usepackage[version=4]{mhchem}

\usepackage{stmaryrd}

\usepackage{graphicx}

\usepackage[export]{adjustbox}

\graphicspath{ {./images/} }



\begin{document}

Коммутаторы операторов (3), действующих на нек-рую скалярную функцию, имеют вид:



\[

\left.\begin{array}{l}

\Delta D-D \Delta=(\gamma+\bar{\gamma}) D+(\varepsilon+\bar{\varepsilon}) \Delta-(\tau+\bar{\pi}) \bar{\delta}-(\bar{\tau}+\pi) \delta, \\

\delta D-D \delta=(\bar{\alpha}+\beta-\bar{\pi}) D+k \Delta-\sigma \bar{\delta}-(\bar{\rho}+\varepsilon-\bar{\varepsilon}) \delta, \\

\delta \Delta-\Delta \delta=-v D+(\tau-\bar{\alpha}-\beta) \Delta+(\mu-\gamma+\bar{\gamma}) \delta, \\

\bar{\delta} \delta-\delta \bar{\delta}=(\bar{\mu}-\mu) D+(\bar{\rho}-\rho) \Delta-(\bar{\alpha}-\beta) \bar{\delta}-(\bar{\beta}-\alpha) \delta .

\end{array}\right\}

\]



Скаляры Ньюмена - Пенроуза. При записи в этих обозначениях различных тензорных соотношений входящие в них тензоры должны быть спроектированы на векторы изотропной тетрады. Исходя из алгебраич. свойств данного тензора, из полного набора его проекщий на векторы изотропной тетрады следует выделить набор алгебраически независимых величин (называемых скаля рам и НьюменаПен роуза), для к-рых приняты специальные обозначения. Приведем здесь соответствующие определения для тензоров Максвелла, Вейля и Риччи.



Те н зор Максвелла электромагнитного поля представляет собой вещественный бивектор (антисимметричный тензор 2-го ранга). Множество бивекторов в 4-мерном пространстве образуют 6-мерное векторное пространство. Шесть линейно независимых бивекторов:



\[

\begin{array}{lll}

U_{i j}=l_{[i} m_{j]}, & V_{i j}=n_{[i} \bar{m}_{j]}, & M_{i j}=l_{[i} n_{j]}-m_{[i} \bar{m}_{j]}, \\

\bar{U}_{i j}=l_{[i} \bar{m}_{j]}, & \bar{V}_{i j}=n_{[i} m_{j]}, & \bar{M}_{i j}=l_{[i} n_{j]}-\bar{m}_{[i} m_{j]},

\end{array}

\]



составляют базис в этом пространстве.



Всякий вещественный бивектор $F_{i j}$ в 4-мерном пространстве относительно так наз. дуального преобразования, имеющего вид:



\[

F_{i j} \propto \stackrel{*}{F_{i j}}, \quad \text { где } \quad \stackrel{*}{F_{i j}}=\frac{1}{2} \varepsilon_{i j k l} F^{k l}

\]



и $\varepsilon_{i j k l}-$ тензор Леви-Чивита $\left(\varepsilon_{i j k l}=\varepsilon_{[i j k l]}, \varepsilon_{0123}=-\sqrt{-g}\right)$, можно разбить на две части - автодуальную $\stackrel{+}{F_{i j}}$, то есть инвариантную относительно дуального преобразования, и антиавтодуальную, к-рая при этом преобразовании умножается на -1 . При этом сам бивектор $F_{i j}$ однозначно характеризуется своей автодуальной частью $\stackrel{+}{F}_{i j}$ :



\[

\stackrel{+}{F_{i j}}=\frac{1}{2}\left(F_{i j}+i \stackrel{*}{F_{i j}}\right) ; \quad F_{i j}=\stackrel{+}{F_{i j}}+\stackrel{\mp}{F_{i j}} .

\]



В то же время легко видеть, что бивекторы (6) автодуальны, тогда как бивекторы (7) антиавтодуальны. Поэтому разложение всякого автодуального бивектора $F_{i j}$ по базису (6), (7) содержит только бивекторы (6), но не (7) [бивекторы (6) образуют базис в 3-мерном комплексном пространстве автодуальных бивекторов]:



\[

\stackrel{+}{F_{i j}}=-2 \varphi_{0} V_{i j}-2 \varphi_{1} M_{i j}+2 \varphi_{2} U_{i j}

\]



Здесь числовые множители в коэффициентах введены для упрощения последующих формул. Таким образом, тензор Максвелла, как и всякий вещественный бивектор, можно полностью охарактеризовать с помощью трех комплексных скаляров Ньюмена - Пенроуза $\varphi_{0}, \varphi_{1}, \varphi_{2}$, определяемых как проекции этого тензора на изотропную тетраду:



\[

\begin{aligned}

& \varphi_{0}=\stackrel{+}{F}_{i j}^{i j}=F_{i j} l^{i} m^{j}, \\

& \varphi_{1}=\frac{1}{2} \stackrel{+}{F}_{i j} M^{i j}=\frac{1}{2} F_{i j}\left(l^{i} n^{j}-m^{i} \bar{m}^{j}\right), \\

& \varphi_{2}=-\stackrel{+}{F_{i j}} V^{i j}=-F_{i j} n^{i} \bar{m}^{j} .

\end{aligned}

\]



Тензор Вейля $C_{i j k l}$ возникает при разбиении тензора кривизны $R_{i j k l}$ на неприводимые части, определяемые бесследовыми тензорами - четырехвалентным $C_{i j k l}$ (тензор Вейля) и двухвалентным $\Phi_{i j}$ и скаляром $\Lambda$ :



\[

\begin{gathered}

R_{i j k l}=C_{i j k l}+\Phi_{i k} g_{j l}-\Phi_{i l} g_{j k}+\Phi_{j l} g_{i k}-\Phi_{j k} g_{i l}+ \\

+2 \Lambda\left(g_{i l} g_{j k}-g_{i k} g_{j l}\right)

\end{gathered}

\]



где $2 \Phi_{i j}=R_{i j}-\frac{1}{4} R g_{i j}-$ бесследовая часть (девиатор) тензора Риччи $R_{i j}\left(R_{i j}=R_{i k j}^{k}\right) ; R-$ скалярная кривизна $\left(R=R_{k}^{k}\right)$; $R=-24 \Lambda$.



Алгебраич. свойства тензора Вейля выражаются соотношениями:



\[

\left.\begin{array}{c}

C_{i j k l}=-C_{j i k l}=-C_{i j k}=C_{k l i j}, \\

C_{i j k l}+C_{i l j k}+C_{i k l j}=0 ; \quad C_{i k j}^{k}=0 .

\end{array}\right\}

\]



Тензор Вейля $C_{i j k l}$ также однозначно описывается своей автодуальной частью согласно формулам:



\[

\stackrel{+}{C}_{i j k l}=\frac{1}{2}\left(C_{i j k l}+i \stackrel{*}{C}_{i j k l}\right) ; \quad C_{i j k l}={\stackrel{+}{C_{i j k l}}}+\stackrel{\mp}{C_{i j k l}},

\]



где $\stackrel{*}{C}_{i j k l}=\frac{1}{2} \varepsilon_{k l m n} C_{i j}{ }^{m n}=\frac{1}{2} \varepsilon_{i j m n} C_{k l}^{m n}$. Так как тензор $\stackrel{+}{C}_{i j k l}$ является автодуальным как по первой, так и по второй паре своих антисимметричных индексов, то он может быть разложен в сумму попарных произведений автодуальных бивекторов (6), то есть представлен в виде симметричной билинейной формы от этих бивекторов. Кроме того, так как в силу (9) имеет место $\stackrel{+}{C}_{i j k}=0$, то среди коэффициентов этой билинейной формы остается лишь пять алгебраически независимых комплексных коэффициентов, обозначаемых $\Psi_{0}, \Psi_{1}, \Psi_{2}, \Psi_{3}, \Psi_{4}$ и определяющих полностью в заданной изотропной тетраде компоненты тензора Вейля:



\[

\begin{gathered}

\stackrel{+}{C}_{i j k l}=-4 \Psi_{0} V_{i j} V_{k l}-4 \Psi_{1}\left(V_{i j} M_{k l}+M_{i j} V_{k l}\right)- \\

\quad-4 \Psi_{2}\left(M_{i j} M_{k l}-U_{i j} V_{k l}-V_{i j} U_{k l}\right)+ \\

\quad+4 \Psi_{3}\left(U_{i j} M_{k l}+M_{i j} U_{k l}\right)-4 \Psi_{4} U_{i j} U_{k l} .

\end{gathered}

\]



Коэффициенты $\Psi_{0}, \Psi_{1}, \Psi_{2}, \Psi_{3}, \Psi_{4}$ вычисляются по формулам:



\[

\begin{aligned}

& \Psi_{0}=-\stackrel{+}{C}_{i j k l} U^{i j} U^{k l}=-C_{i j k l} l^{i} m^{j} l^{k} m^{l} \\

& \Psi_{1}=-\frac{1}{2} \stackrel{+}{C}_{i j k l} U^{i j} M^{k l}=-C_{i j k l} l^{i} n^{j} l^{k} m^{l}

\end{aligned}

\]



\begin{center}

\includegraphics[max width=\textwidth]{2023_07_08_48d3224d4755f11871c2g-1}

\end{center}



\[

\begin{aligned}

& \Psi_{3}=\frac{1}{2} \stackrel{+}{C}_{i j k l} V^{i j} M^{k l}=-C_{i j k l} n^{i} l^{j} n^{k} \bar{m}^{l}, \\

& \Psi_{4}=-\stackrel{+}{C}_{i j k l} V^{i j} V^{k l}=-C_{i j k l} n^{i} \bar{m}^{j} n^{k} \bar{m}^{l} .

\end{aligned}

\]



Для проекций тензора $\Phi_{i j}-$ бесследовой части те н з о р а Р и ч ч и приняты следующие обозначения:



\[

\begin{aligned}

& \Phi_{00}=\Phi_{i j} l^{i} l^{j}, \quad \Phi_{10}=\Phi_{i j} l^{i} \bar{m}^{j}, \quad \Phi_{20}=\Phi_{i j} \bar{m}^{i} \bar{m}^{j}, \\

& \Phi_{01}=\Phi_{i j} l^{i} m^{j}, \quad \Phi_{11}=\Phi_{i j} l^{i} n^{j}, \quad \Phi_{21}=\Phi_{i j} n^{i} \bar{m}^{j}, \\

& \Phi_{02}=\Phi_{i j} m^{i} m^{j}, \quad \Phi_{12}=\Phi_{i j} n^{i} m^{j}, \quad \Phi_{22}=\Phi_{i j} n^{i} n^{j} .

\end{aligned}

\]



При этом в силу симметрии и бесследовости $\Phi_{i j}$ для его проекций (11) имеют место соотношения: $\Phi_{m n}=\bar{\Phi}_{n m}$, где $m, n=0,1,2$. Поэтому только шесть величин из (11) (три вещественных и три комплексных) независимы.



Из сказанного выше вытекает, что пять комплексных величин (10), шесть независимых скаляров (11) и скаляр $\Lambda$ определяют полный набор проекций на изотропную тетраду тензора Римана-Кристоффеля $R_{i j k l}$.



Уравнения гравитационного поля в формализме Ньюмена Пенроуза. В обычной формулировке общей теории относительности гравитационное поле описывается компонента-





\end{document}